\section*{Quick start}
\addcontentsline{toc}{section}{Quick start}

\noindent\textbf{From zero to the Weisshorn results in four commands} (Windows; see Section~\ref{sec:install}
for macOS/Linux):
\begin{lstlisting}[language=bash]
python -m venv .venv
.venv\Scripts\activate
pip install -r requirements.txt
python -m nosecone run configs/weisshorn.yaml
\end{lstlisting}
\note{Then open \texttt{output/weisshorn/summary.md}: it lists every nose of the study, best first, with its
wave drag, apogee, mass, payload room and static-margin change. Figures are in
\texttt{output/weisshorn/figures/}, CAD files (CSV, STL, STEP) in \texttt{output/weisshorn/cad/}.}

\medskip
\noindent\textbf{To study another rocket}, copy \texttt{configs/example\_new\_rocket.yaml}, change the values of
your vehicle (diameter, length, masses, engine, fins) and the noses you want to compare, and run it with the same
command. You never need to edit the Python code for that: every input is in the YAML file. The parameters
are described in Section~\ref{sec:parameters}, the outputs in Section~\ref{sec:outputs}.

\medskip
\noindent\textbf{How this guide is organised.}
\begin{itemize}[leftmargin=*,itemsep=1pt]
\item Section~\ref{sec:scope}: what the tool computes, with which models, and where these models stop being valid.
\item Section~\ref{sec:install}: installation and requirements.
\item Section~\ref{sec:structure}: the repository and what each module does.
\item Section~\ref{sec:parameters}: the configuration file and the parameters you change.
\item Section~\ref{sec:running}: running a single design and a parameter sweep, with worked examples for our 180~mm rocket.
\item Section~\ref{sec:outputs}: how to read the plots and tables.
\item Section~\ref{sec:extending}: adding a profile family, a plot, a parameter, or changing a model.
\end{itemize}
